\par\bigskip
\SourceLeaf{489}%
\pdfbookmark[0]{Plan for the degree survey. 1819.}{entwurf-gradmessung-1819}
\begin{center}\large [PLAN FOR THE DEGREE SURVEY.]\end{center}
\par
— — —

I have learned from a letter from Professor \textsc{Schumacher} that during the professor’s stay in London, the Danish envoy there, Count \textsc{Bourke}, spoke with Count \textsc{Münster} about the benefit of continuing the Danish degree survey through the Kingdom of Hanover, and that the latter asked me to write to him on the subject. I felt bound to regard this communication as an order. The warm interest Your — — have taken in this scientific matter, and shown as early as last year by taking measures to secure the possibility of such a continuation, obliges me to inform Your — — immediately of this circumstance and assures me that you will kindly receive my request to support the matter in London as far as possible. I have all the less reason to hesitate to make this request, since I am acting solely out of love for science. Without that love and the lively conviction of the great importance of such an undertaking, I would be more inclined, given my personal preferences, to regard a prolonged absence from purely astronomical work as a sacrifice. As Professor \textsc{Schumacher} writes to me, the venerable Sir \textsc{Joseph Banks} too has taken up with great warmth the idea of extending this degree survey farther on the continent and has promised his full cooperation.

Professor \textsc{Schumacher} will himself have told Your — — as he passed through Hanover that the purpose of his journey to England was to take delivery of the famous large zenith sector used in the English degree survey, which is being lent to him for his astronomical observations. Immediately after his return he will begin these observations at Lauenburg, the southern endpoint of his survey. Given the prospect of continuing the Danish degree survey through the Kingdom of Hanover, it would undoubtedly be very desirable for both surveys in many respects if I made \textit{these} observations jointly with him. \textit{Without} this cooperation, it is easy to imagine discrepancies or doubts arising in future that would then be difficult to clear up and that would be forestalled \textit{by} it. Concerns did in fact arise later precisely about the astronomical observations made at Dunkirk, the northern endpoint of the French degree survey. On that account it was found necessary for English and French commissioners to make new observations jointly with the zenith sector in question last autumn. And I hear that as soon as Professor \textsc{Schumacher} has returned it, this very zenith sector is to be shipped to the Balearic Islands for yet another repetition of the observations there. Only by observing with this instrument myself and comparing my experience with what I have learned using \textsc{Repsold}’s meridian circle, and still more with what I will learn using \textsc{Reichenbach}’s circle, would I be able to judge whether it will be necessary or especially desirable to bring in that very zenith sector for the observations in the Kingdom of Hanover. Even if the continuation of the degree survey were never carried out, the experience I would gain with the zenith sector, an instrument unique of its kind, would not be without substantial benefit for my future observations with the meridian instruments of the observatory here. As evidence of how important the observations in Lauenburg are considered, I may add that Dr. \textsc{Olbers} has written to me of his intention to come to Lauenburg to familiarize himself with the instrument and attend the observations.
